\section{Preprocessing: tiling antigens into assay constructs}\label{sec:tiling}
The no-antibody (DP4) work starts before any design is generated, by turning an antigen into
a set of short peptide constructs to scaffold. Two things drive how we do this.

First, RFD3 needs an input crystal structure. The contig tells RFD3 which residues to hold
fixed while it builds a scaffold around them, so those residues have to exist in a PDB. We
therefore tile the \emph{PDB-resolved} sequence (the residues actually present in the
structure), not the deposited FASTA. 6M0J is the clearest case: its FASTA starts 14 residues
before the crystal (chain E begins at residue 333, which is position 15 in the FASTA), so we
drop those 14 and tile the 223-residue resolved region. 1D2K (392) and 4WAT (402) match
their full chains.

Second, the constructs have to match the linear-epitope assay. Each tile is placed at the
C-terminus of a constant 73-residue construct: a glycine--serine filler followed by the
ENLYFQGA protease site. Because the prefix is constant, every member is synthesized at a
constant length (103 residues for a 30-mer), and the protease site cleaves it down to the
bare tile at the end, which is what the assay reads.

\subsection{General tiling utility}
The tiling step is a single utility (\texttt{episcaf\_pipeline/tile\_fasta.py}) so we can
point it at any target. It takes a sequence source (a FASTA, or the \texttt{antigen\_seq}
column of a ledger such as \texttt{dp2.parquet}), a window length, and a step, and writes the
assay library in the DP2 annotated format (\texttt{library\_member}, \texttt{sequence},
\texttt{category}, \texttt{model}, \texttt{designedSequence}, \texttt{designedSequenceLength},
\texttt{design\_ID}, \texttt{target}). Windows are placed at the chosen step; when the last
window stops short of the C-terminus, we add one final window so the C-terminal residues are
still covered. We pinned the utility to what was already run: re-tiling 1D2K, 4WAT, and 6M0J
at a 12-residue window and step 2 reproduces the existing 12-mer library exactly (191, 196,
and 107 tiles, every tile and full construct identical).

\subsection{Tiled 12-mers, and tiled 30-mers as linear controls}
The existing set tiles three antigens into 12-mers. The next category is tiled 30-mers,
which serve as the \emph{linear controls}: unscaffolded constructs (model \texttt{(none)})
that present the bare peptide, exactly what the scaffolded designs are meant to beat. For
those we tile all 59 mAb-target antigens (sequences taken directly from \texttt{dp2.parquet})
plus the three tiled antigens, at a 30-residue window and step 6.

\subsection{Why the design count is not (tiles $\times$ top-$k$)}
A natural expectation is that taking the top $k$ designs per epitope gives $k$ times the
number of tiles. It does not, and the gap is informative. Each tile is one epitope that RFD3
scaffolds into many designs (seeds $\times$ ProteinMPNN sequences), which are then
AlphaFold3-validated and gated; we keep at most $k$ per epitope, so the selected total is
$\sum_{\text{epitopes}} \min(k, n_{\text{gated}})$. In practice the per-epitope yield is not
the bottleneck: every epitope that survives has well over $k$ designs (median $\sim$190
scored). The shortfall is whole epitopes that drop out, because a tile whose residues are
unresolved in the crystal cannot be aligned to the native antigen and is discarded at scoring
(\texttt{crystal\_epi\_incomplete}). Table~\ref{tab:tilecount} shows this for top-5 selection.

\begin{table}[h]\centering
\begin{tabular}{lcccc}
\toprule
antigen & tiles & epitopes lost (crystal) & epitopes kept & top-5 designs \\
\midrule
1D2K & 191 & 0  & 191 & 951 \\
6M0J & 107 & 15 & 92  & 450 \\
4WAT & 196 & 40 & 156 & 780 \\
\bottomrule
\end{tabular}
\caption{Selected designs track (epitopes kept)$\,\times\,5$, not (tiles)$\,\times\,5$.
1D2K loses no epitopes; 6M0J and 4WAT lose 15 and 40 tiles to unresolved crystal regions.
The small residuals (951 vs 955, 450 vs 460) are epitopes with fewer than 5 designs clearing
the \SI{2.5}{\angstrom} gate.}
\label{tab:tilecount}
\end{table}
